%				standard font size
%command		10pt	11pt	12pt

%\tiny			5pt		6pt		6pt
%\scriptsize	7pt		8pt		8pt
%\footnotesize	8pt		9pt		10pt
%\small			9pt		10pt	11pt
%\normalsize	10pt	11pt	12pt
%\large			12pt	12pt	14pt
%\Large			14pt	14pt	17pt
%\LARGE			17pt	17pt	20pt
%\huge			20pt	20pt	25pt
%\Huge			25pt	25pt	25pt
%								 ||	

%======================================================
% DECLARACIÓN DE PAQUETES A UTILIZAR
% (se pueden agregar a la lista cuantos paquetes sean necesario declarar )
%======================================================
% babel[spanish] requires texlive-lang-spanish which may not be installed.
% Using english babel for framework compatibility; accented chars handled via inputenc+T1.
\usepackage[english]{babel}
% \selectlanguage{spanish}  % commented out — not available without lang pack
\usepackage[utf8]{inputenc} % latin1 es para sistemas basados en windows, para sistemas unix la codificación de caracteres es utf8
\usepackage[T1]{fontenc} %paquete necesario para los acentos, NO es necesario usar (\'a) para obtener (á)
\usepackage{amssymb,amsmath,amsfonts}
\usepackage{color,curves}
\usepackage{epsfig}
\usepackage{times}
\usepackage{multirow,rotating}
\usepackage{times}
\usepackage{subfig,wrapfig}
\usepackage{pdfpages}
\usepackage[bookmarks=true,citebordercolor={1 1 1},linkbordercolor={1 1 1},urlbordercolor={1 1 1},hidelinks]{hyperref}
\usepackage{bookmark}
\usepackage{booktabs}
\usepackage{lipsum}
\usepackage{longtable}
\usepackage{adjustbox}
\usepackage{array}
\usepackage{enumitem}
\usepackage{colortbl, xcolor}
\usepackage{bm}	
\usepackage{authoraftertitle}
\usepackage{mfirstuc}
\MFUnocap{de}
\MFUnocap{y}
\MFUnocap{por}
\MFUnocap{para}
\MFUnocap{del}
\MFUnocap{un}
\MFUnocap{una}
\MFUnocap{el}
\MFUnocap{la}
\MFUnocap{en}
\MFUnocap{con}
% \usepackage[a4paper,left=3cm,right=3cm,top=2.5cm,bottom=2.5cm,showframe]{geometry}
\usepackage[a4paper,left=3cm,right=3cm,top=2.5cm,bottom=2.5cm]{geometry}
\usepackage{longtable}
\usepackage{tabularx}
\usepackage{tabulary}
%\usepackage{cite}
\usepackage[notlof, notlot, numbib,nottoc]{tocbibind}
%\usepackage[none]{hyphenat}
\usepackage{soul}
\usepackage{anyfontsize}
\usepackage{capt-of}
\usepackage{ifthen}
\usepackage{float}
\usepackage{pdflscape}
\usepackage{xurl}
\usepackage{import}
\usepackage{graphicx}
\usepackage{footnote}
\usepackage{changepage}
%\usepackage{showframe}
\usepackage{pgfgantt}
\usepackage{nameref}
%\renewcommand\ShowFrameLinethickness{0.25pt}
%\renewcommand*\ShowFrameColor{\color{green}}
\usepackage{helvet}
\renewcommand{\familydefault}{\sfdefault}
% IEEEtrantools: using local stub (full package may not be installed)
\usepackage{IEEEtrantools}
\usepackage{easyReview} % for \comment, \alert, etc.
\usepackage{listings}
\lstset{
    language=bash, %% Troque para PHP, C, Java, etc... bash é o padrão
    basicstyle=\ttfamily\small,
    numberstyle=\footnotesize,
    numbers=left,
    backgroundcolor=\color{gray!10},
    frame=single,
    tabsize=2,
    rulecolor=\color{black!30},
    title=\lstname,
    escapeinside={\%*}{*)},
    breaklines=true,
    breakatwhitespace=true,
    framextopmargin=2pt,
    framexbottommargin=2pt,
    inputencoding=utf8,
    showstringspaces=false,
    extendedchars=true,
    literate={á}{{\'a}}1 {ñ}{{\~n}}1 {é}{{\'e}}1 {ó}{{\'o}}1 {í}{{\'i}}1
          {á}{{\'a}}1  {é}{{\'e}}1  {í}{{\'i}}1 {ó}{{\'o}}1  {ú}{{\'u}}1
      {Á}{{\'A}}1  {É}{{\'E}}1  {Í}{{\'I}}1 {Ó}{{\'O}}1  {Ú}{{\'U}}1
      {à}{{\`a}}1  {è}{{\`e}}1  {ì}{{\`i}}1 {ò}{{\`o}}1  {ù}{{\`u}}1
      {À}{{\`A}}1  {È}{{\`E}}1  {Ì}{{\`I}}1 {Ò}{{\`O}}1  {Ù}{{\`U}}1
      {ä}{{\"a}}1  {ë}{{\"e}}1  {ï}{{\"i}}1 {ö}{{\"o}}1  {ü}{{\"u}}1
      {Ä}{{\"A}}1  {Ë}{{\"E}}1  {Ï}{{\"I}}1 {Ö}{{\"O}}1  {Ü}{{\"U}}1
      {â}{{\^a}}1  {ê}{{\^e}}1  {î}{{\^i}}1 {ô}{{\^o}}1  {û}{{\^u}}1
      {Â}{{\^A}}1  {Ê}{{\^E}}1  {Î}{{\^I}}1 {Ô}{{\^O}}1  {Û}{{\^U}}1
      {œ}{{\oe}}1  {Œ}{{\OE}}1  {æ}{{\ae}}1 {Æ}{{\AE}}1  {ß}{{\ss}}1
      {ẞ}{{\SS}}1  {ç}{{\c{c}}}1 {Ç}{{\c{C}}}1 {ø}{{\o}}1  {Ø}{{\O}}1
      {å}{{\aa}}1  {Å}{{\AA}}1  {ã}{{\~a}}1  {õ}{{\~o}}1 {Ã}{{\~A}}1
      {Õ}{{\~O}}1  {ñ}{{\~n}}1  {Ñ}{{\~N}}1  {¿}{{?`}}1  {¡}{{!`}}1
      {°}{{\textdegree}}1 {º}{{\textordmasculine}}1 {ª}{{\textordfeminine}}1
      {£}{{\pounds}}1  {©}{{\copyright}}1  {®}{{\textregistered}}1
      {«}{{\guillemotleft}}1  {»}{{\guillemotright}}1  {Ð}{{\DH}}1  {ð}{{\dh}}1
      {Ý}{{\'Y}}1    {ý}{{\'y}}1    {Þ}{{\TH}}1    {þ}{{\th}}1    {Ă}{{\u{A}}}1
      {ă}{{\u{a}}}1  {Ą}{{\k{A}}}1  {ą}{{\k{a}}}1  {Ć}{{\'C}}1    {ć}{{\'c}}1
      {Č}{{\v{C}}}1  {č}{{\v{c}}}1  {Ď}{{\v{D}}}1  {ď}{{\v{d}}}1  {Đ}{{\DJ}}1
      {đ}{{\dj}}1    {Ė}{{\.{E}}}1  {ė}{{\.{e}}}1  {Ę}{{\k{E}}}1  {ę}{{\k{e}}}1
      {Ě}{{\v{E}}}1  {ě}{{\v{e}}}1  {Ğ}{{\u{G}}}1  {ğ}{{\u{g}}}1  {Ĩ}{{\~I}}1
      {ĩ}{{\~\i}}1   {Į}{{\k{I}}}1  {į}{{\k{i}}}1  {İ}{{\.{I}}}1  {ı}{{\i}}1
      {Ĺ}{{\'L}}1    {ĺ}{{\'l}}1    {Ľ}{{\v{L}}}1  {ľ}{{\v{l}}}1  {Ł}{{\L{}}}1
      {ł}{{\l{}}}1   {Ń}{{\'N}}1    {ń}{{\'n}}1    {Ň}{{\v{N}}}1  {ň}{{\v{n}}}1
      {Ő}{{\H{O}}}1  {ő}{{\H{o}}}1  {Ŕ}{{\'{R}}}1  {ŕ}{{\'{r}}}1  {Ř}{{\v{R}}}1
      {ř}{{\v{r}}}1  {Ś}{{\'S}}1    {ś}{{\'s}}1    {Ş}{{\c{S}}}1  {ş}{{\c{s}}}1
      {Š}{{\v{S}}}1  {š}{{\v{s}}}1  {Ť}{{\v{T}}}1  {ť}{{\v{t}}}1  {Ũ}{{\~U}}1
      {ũ}{{\~u}}1    {Ū}{{\={U}}}1  {ū}{{\={u}}}1  {Ů}{{\r{U}}}1  {ů}{{\r{u}}}1
      {Ű}{{\H{U}}}1  {ű}{{\H{u}}}1  {Ų}{{\k{U}}}1  {ų}{{\k{u}}}1  {Ź}{{\'Z}}1
      {ź}{{\'z}}1    {Ż}{{\.Z}}1    {ż}{{\.z}}1    {Ž}{{\v{Z}}}1,
}

\sodef\spaceout{}{0pt plus 1fil}{.8em plus 1fil}{0pt}
\definecolor{upiizcolor}{RGB}{113,48,70}

% Colores y macros para cronogramas Gantt (Anexo1 y Anexo2)
\definecolor{myyellow}{HTML}{FFC000}
\definecolor{mygreen}{HTML}{00B050}
\definecolor{myblue}{HTML}{00B0F0}
\definecolor{mygray}{rgb}{0.816,0.816,0.816}
\definecolor{myred}{rgb}{0.851,0,0}
\newcommand{\cg}{\cellcolor{mygreen}}
\newcommand{\cb}{\cellcolor{myblue}}
\newcommand{\cy}{\cellcolor{myyellow}}
\newcommand{\cgr}{\cellcolor{mygray}}
\newcommand{\cred}{\cellcolor{myred}}
\newcommand{\ck}{\cellcolor{black!85}}
\newcolumntype{C}{>{\centering\arraybackslash}p{0.35cm}}

\usepackage{datatool}
\newcommand{\sortitem}[1]{
	\DTLnewrow{list}
	\DTLnewdbentry{list}{description}{#1}
}
\newenvironment{sortedlist}{
	\DTLifdbexists{list}{\DTLcleardb{list}}{\DTLnewdb{list}}
}{%
	\DTLsort{description}{list}
	\begin{itemize}
		\DTLforeach*{list}{\theDesc=description}{
			\item \theDesc}
	\end{itemize}
}

\newcommand{\itemcolor}[1]{
	\renewcommand{\makelabel}[1]{\color{#1}\hfil ##1}}
\setlist[itemize]{label={$\diamond$}}
\setlength\parindent{0pt}

\counterwithout{footnote}{chapter}

%======================================================
% DECLARACIÓN DE CARPETA DONDE SE ENCUENTRAN LAS FIGURAS
%======================================================

% NOTA MUY IMPORTANTE

% Dependiendo del formato de las figuras será el método de compilación del documento:

%*.EPS o *.PDF ---- Latex=>dvi2ps=>ps2pdf=>ver pdf
%					Latex=>dvi2pdf=>ver pdf
%*.png y *.jpg ---- PDFLatex=>ver PDF					

\graphicspath{{Figuras/}}  % Las figuras se encuentran en una carpeta llamada 'Figuras' que se incluye en la carpeta principal de los archivos .tex
%================================================================
% OPCIONES DE DOCUMENTO
%
% La opción \psdraft deja un espacio en blanco en donde van
% las figuras. Esto ahorra toner en la impresión del borrador de tesis
%
% La opción \psfull hace incluir a todas las figuras
%
%===============================================================

\psfull

%===============================================================
% ESTILO DE PÁGINA
%
% El estilo de página \pagestyle{thesis} coloca los encabezados
% correctos de la versión final del documento
%==============================================================
\pagestyle{thesis}

%==============================================================
\noappendixtables
\noappendixfigures

\makeatletter
\let\@openbib@code\relax
\makeatother

\makeatletter
\let\newtitle\@title
\let\newauthor\@author
\let\newauthortitle\@authortitle
\let\newdate\@date
\makeatother

\providecommand{\decimalpoint}{}

\counterwithout{footnote}{chapter}